\subsection{THE COMPACT-OPEN TOPOLOGY}

THEOREM 2. Let X be a topological space, Y a uniform space. For each compact subset K of X and each open subset U of Y, let \(\mathbf{T}(\mathbf{K},\mathbf{U})\) denote the set of all continuous mappings \(u:X\to Y\) such that \(u(\mathbf{K})\subset U\) . Then the sets \(\mathbf{T}(\mathbf{K},\mathbf{U})\) generate the topology of compact convergence on \(\mathcal{C}(\mathbf{X};\mathbf{Y})\) .

Let \(Y'\) be the Hausdorff uniform space associated with \(Y\) (Chapter II, § 3, no. 8) and let \(i: Y \to Y'\) be the canonical mapping of \(Y\) onto \(Y'\) . The topology of compact convergence is the coarsest topology for which the mappings \(u \to (i \circ u)|K\) of \(\mathcal{C}(X; Y)\) into \(\mathcal{C}_u(K; Y')\) are continuous, as \(K\) runs through the set of all compact subsets of \(X\) (§ 1, no. 4, Proposition 4). Hence we obtain a subbase of the topology of \(\mathcal{C}_c(X; Y)\) by taking a subbase of the topology of \(\mathcal{C}_c(K; Y')\) for each compact subset \(K\) of \(X\) and then taking the union {[}in \(\mathfrak{P}(\mathcal{C}(X; Y))\) {]} of the inverse images of these subbases in \(\mathcal{C}(X, Y)\) . On the other hand, every open subset of \(Y\) is of the form \(\overline{i}^1(U')\) , where \(U'\) is open in \(Y'\) (Chapter II, § 3, no. 7, Proposition 12); hence, for each compact subset \(K' \supset K\) , \(T(K, \overline{i}^1(U'))\) is the inverse image of \(T(K, U')\) under the mapping

\[
\mathcal {C} (\mathbf {X}; \mathbf {Y}) \rightarrow \mathcal {C} _ {u} (\mathbf {K} ^ {\prime}; \mathbf {Y} ^ {\prime}).
\]

It thus remains for us to prove the theorem when X is compact and Y is Hausdorff; we shall make these assumptions from now on.

Let us first show that \(T(\mathbf{K}, \mathbf{U})\) is open in \(\mathcal{C}_c(\mathbf{X}; \mathbf{Y})\) . Let \(u_0\) be a point of this set; since \(u_0(\mathbf{K})\) is compact (Chapter I, § 9, no. 4, Theorem 2, Corollary 1) and contained in the open set \(\mathbf{U}\) , there exists a symmetric entourage \(\mathbf{V}\) of \(\mathbf{Y}\) such that \(\mathbf{V}(u_0(\mathbf{K})) \subset \mathbf{U}\) (Chapter II, § 4, no. 3, Proposition 4, Corollary). Let \(\mathbf{W}\) be the neighbourhood of \(u_0\) in \(\mathcal{C}_c(\mathbf{X}; \mathbf{Y})\) consisting of all continuous mappings \(u: \mathbf{X} \to \mathbf{Y}\) such that \((u(x), u_0(x)) \in \mathbf{V}\) for all \(x \in \mathbf{K}\) . For such mappings we clearly have \(u(\mathbf{K}) \subset \mathbf{V}(u_0(\mathbf{K})) \subset \mathbf{U}\) ; hence \(u \in T(\mathbf{K}, \mathbf{U})\) and therefore \(W \subset T(\mathbf{K}, \mathbf{U})\) , which proves our assertion.

Conversely, if W is a neighbourhood of a point \(u_{0} \in \mathcal{C}_{c}(\mathbf{X}; \mathbf{Y})\) , let us show that W contains the intersection of a finite number of neighbourhoods of the form \(T(\mathbf{K}, \mathbf{U})\) . We may suppose that W is the set of all \(u \in \mathcal{C}(\mathbf{X}; \mathbf{Y})\) such that \((u(x), u_{0}(x)) \in \mathbf{V}\) for all \(x \in X\) , V being a given entourage of Y. Since \(u_{0}\) is continuous on X, it is uniformly continuous (Chapter II, § 4, no. 1, Theorem 2). Let \(V_{1}\) be a symmetric entourage of Y, open in \(Y \times Y\) and such that \(V_{1}^{2} \subset V\) . X can be covered by a finite number of compact sets \(K_{i}\) ( \(1 \leqslant i \leqslant n\) ) such that each \(u_{0}(\mathbf{K}_{i})\) is \(V_{1}\) -small ( \(1 \leqslant i \leqslant n\) ). Let \(U_{i}\) be the open set \(V_{1}(u_{0}(\mathbf{K}_{i}))\) , and let \(u: X \to Y\) be a continuous mapping contained in the intersection of the n sets \(T(\mathbf{K}_{i}, \mathbf{U}_{i})\) (which are neighbourhoods of \(u_{0}\) ). Then, for every \(x \in K_{i}\) , \(u(x)\) belongs to \(U_{i}\) and therefore \(u_{0}(x)\) and \(u(x)\) are \(V_{1}^{2}\) -close, hence V-close. Since each \(x \in X\) belongs to some \(K_{i}\) , we have \(u \in W\) and the proof is complete.

This result leads us to make the following definition :

DEFINITION 1. Let X, Y be two topological spaces, not necessarily uniformizable. For each compact subset K of X and each open subset U of Y, let T(K, U) be the set of all \(u \in \mathcal{C}(X; Y)\) such that \(u(K) \subset U\) . The topology on C(X; Y) generated by the sets T(K, U) is called the topology of compact convergence or the compact-open topology; and we denote by \(\mathcal{C}_{c}(X; Y)\) the topological space obtained by endowing C(X; Y) with this topology.

If Y is a uniform space it follows from Theorem 2 that this definition agrees with that given in §1, no. 3.

If H is a subset of \(\mathcal{C}(\mathbf{X};\mathbf{Y})\) we shall say that the topology induced on H by that of \(\mathcal{C}_{c}(\mathbf{X};\mathbf{Y})\) is the compact-open topology on H.

Example. Let I be the interval \([0, 1]\) in R. If Y is any topological space, the space \(\mathcal{C}_{c}(I; Y)\) is called the space of paths in Y. For each \(y \in Y\) , the subspace \(\Omega_{y}(Y)\) of \(\mathcal{C}_{c}(I; Y)\) consisting of paths u such that \(u(0) = u(1) = y\) is called the space of loops (in Y) at the point y.

Remarks. 1) Likewise, the topology induced on \(\mathcal{C}(\mathbf{X};\mathbf{Y})\) by the product topology on \(\mathbf{Y}^{\mathbf{X}} = \mathcal{F}(\mathbf{X};\mathbf{Y})\) is called the topology of pointwise convergence (Y being not necessarily uniformizable); it is generated by sets of the form \(T(\{x\}, \mathbf{U})\) as \(x\) runs through \(\mathbf{X}\) and \(\mathbf{U}\) runs through the set of all open subsets of \(\mathbf{Y}\) , and it is therefore coarser than the compact-open topology. We deduce that, if \(\mathbf{Y}\) is Hausdorff, the space \(\mathcal{C}_s(\mathbf{X}; \mathbf{Y})\) is Hausdorff (Chapter I, § 8, no. 1, Proposition 5, Corollary).

\begin{enumerate}
\def\labelenumi{\arabic{enumi})}
\setcounter{enumi}{1}
\tightlist
\item
  Let \(\mathfrak{S}\) be a subbase of the topology of Y, and let \(\mathfrak{R}\) be a set of compact subsets of X with the following property:
\end{enumerate}

\begin{enumerate}
\def\labelenumi{(\Alph{enumi})}
\setcounter{enumi}{17}
\tightlist
\item
  If \(\mathbf{L}\) is any compact subset of \(\mathbf{X}\) and \(\mathbf{V}\) is any neighbourhood of \(\mathbf{L}\) , there exists a finite number of sets \(\mathbf{K}_i \in \mathfrak{R}\) such that \(\mathbf{L} \subset \bigcup_{i} \mathbf{K}_i \subset \mathbf{V}\) .
\end{enumerate}

Then the sets \(T(\mathbf{K}, \mathbf{U})\) , where \(\mathbf{K} \in \mathfrak{R}\) and \(\mathbf{U} \in \mathfrak{S}\) , form a subbase for the compact-open topology on \(\mathcal{C}(\mathbf{X}; \mathbf{Y})\) . To prove this, we have to show that if \(\mathbf{L}\) is any compact subset of \(\mathbf{X}\) and \(\mathbf{V}\) any open subset of \(\mathbf{Y}\) , and if \(u \in T(\mathbf{L}, \mathbf{V})\) , then there exists a finite number of pairs \((\mathbf{K}_i, \mathbf{U}_i)\) such that \(\mathbf{K}_i \in \mathfrak{R}\) , \(\mathbf{U}_i \in \mathfrak{S}\) and \(u \in \bigcap_{i} T(\mathbf{K}_i, \mathbf{U}_i) \subset T(\mathbf{L}, \mathbf{V})\) . Note first that for every finite sequence \((s_k)\) of sets of \(\mathfrak{S}\) and every compact subset \(\mathbf{M}\) of \(\mathbf{X}\) , we have \(T\left(\mathbf{M}, \bigcap_k \mathbf{S}_k\right) = \bigcap_k T(\mathbf{M}, \mathbf{S}_k)\) by definition. We may therefore first of all replace \(\mathfrak{S}\) by the set of finite intersections of sets of \(\mathfrak{S}\) , i.e.~we may suppose that \(\mathfrak{S}\) is a base of the topology of \(\mathbf{Y}\) . By hypothesis, \(u(\mathbf{L})\) is quasi-compact and contained in \(\mathbf{V}\) , hence there exists a finite number of sets \(\mathbf{U}_i \in \mathfrak{S}\) contained in \(\mathbf{V}\) which cover \(u(\mathbf{L})\) . The sets \(\overline{u}^1(\mathbf{U}_i)\) are open in \(\mathbf{X}\) and cover \(\mathbf{L}\) . For each \(x \in \mathbf{L}\) there is therefore a compact neighbourhood \(N_x\) of \(x\) in \(\mathbf{L}\) , contained in some one of the \(\overline{u}^1(\mathbf{U}_i)\) . We can cover \(\mathbf{L}\) with a finite number of these sets \(N_{x_j} = L_j\) ; for each \(j\) , let us denote by \(i(j)\) one of the indices \(i\) such that \(L_j \subset \overline{u}^1(\mathbf{U}_i)\) . This being so, for each index \(j\) there exists {[}by (R){]} a finite number of sets \(K_{jk} \subset \overline{u}^1(\mathbf{U}_{i(j)})\) , belonging to \(\mathfrak{R}\) , which cover \(L_j\) . For each \(v \in \bigcap_{j,k} T(\mathrm{K}_{jk}, \mathrm{U}_{i(j)})\) we have \(\bigcup_{k} v(\mathrm{K}_{jk}) \subset U_{i(j)}\) and therefore \(v(\mathrm{L}_j) \subset U_{i(j)}\) , and \(v(\mathrm{L}) = \bigcup_{j} v(\mathrm{L}_j) \subset \bigcup_{j} U_{i(j)} \subset V\) ; thus our assertion is proved.

THEOREM 3. Let X, Y, Z be three topological spaces and let f be a mapping of X × Y into Z. If f is continuous then \(f: x \to f(x, .)\) is a continuous mapping of X into \(\mathcal{C}_{c}(Y; Z)\) . The converse is true if Y is locally compact.

Suppose that \(f\) is continuous. To show that \(\tilde{f}\) is continuous we have to prove that, for each compact subset \(K\) of \(Y\) and each open subset \(U\) of \(Z\) , the inverse image \(V\) of \(T(K, U)\) under \(\tilde{f}\) is open in \(X\) . Let \(x_0 \in V\) ; for each \(y \in K\) , we have \(f(x_0, y) \in U\) , and since \(f\) is continuous there is a neighbourhood \(V_y\) of \(x_0\) in \(X\) and a neighbourhood \(W_y\) of \(y\) in \(Y\) such that \(f(V_y \times W_y) \subset U\) . Since \(K\) is compact, there exists a finite number of points \(y_{i} \in K\) such that the sets \(W_{y_{i}}(1 \leqslant i \leqslant n)\) cover K. Let \(V'\) be the intersection of the neighbourhoods \(V_{y_{i}}\) of \(x_{0}\) , which is a neighbourhood of \(x_{0}\) ; if \(x \in V'\) and \(y \in K\) , we have \(f(x, y) \in U\) , since y is contained in some one of the \(W_{y_{i}}\) and x is contained in each \(V_{y_{i}}\) ; hence \(V' \subset V\) , and therefore V is a neighbourhood of each of its points and consequently is open in X.

Conversely, suppose that \(\tilde{f}\) is continuous and that Y is locally compact, and let us show that \(f\) is continuous. Let \(x_0 \in \mathbf{X}\) , let \(y_0 \in \mathbf{Y}\) and let U be an open neighbourhood of \(f(x_0, y_0)\) in Z; we shall show that there is a neighbourhood V of \(x_0\) in X and a neighbourhood W of \(y_0\) in Y such that \(f(\mathbf{V} \times \mathbf{W}) \subset \mathbf{U}\) . Since \(y \to f(x_0, y)\) is continuous, there is a compact neighbourhood W of \(y_0\) such that \(f(\{x_0\} \times \mathbf{W}) \subset \mathbf{U}\) . On the other hand, since \(\tilde{f}\) is continuous, the set V of \(x \in \mathbf{X}\) such that \(f(x, .) \in T(\mathbf{W}, \mathbf{U})\) {[}i.e.~such that \(f(x, y) \in \mathbf{U}\) for all \(y \in \mathbf{W}]\) is an open subset of X and therefore a neighbourhood of \(x_0\) ; and we have \(f(\mathbf{V} \times \mathbf{W}) \subset \mathbf{U}\) .

Q.E.D.

COROLLARY I. Let X be a locally compact space, Y a topological space, H a subset of C (X; Y). Then the compact-open topology on H is the coarsest for which the mapping \((u, x) \to u(x)\) of \(H \times X\) into Y is continuous.

For, by Theorem 3, this mapping is continuous if and only if the canonical injection \(\mathbf{H} \to \mathcal{C}_c(\mathbf{X}; \mathbf{Y})\) is continuous.

Remark. 3) Let X be a locally compact space and Y a Hausdorff topological space. If \(\mathcal{G}\) is a topology on a subset H of \(\mathcal{C}(\mathbf{X};\mathbf{Y})\) such that the mapping \((u,x)\to u(x)\) is continuous on \(\mathrm{H}\times \mathrm{X}\) and if also H is compact with respect to \(\mathcal{G}\) , then \(\mathcal{G}\) is the compact-open topology. For it is finer than the latter by Corollary 1, and since the compact-open topology is Hausdorff, the two topologies are identical. Note that if in addition Y is completely regular, then H is equicontinuous with respect to every uniformity compatible with the topology of Y (§2, no. 5, Theorem 2, Corollary 3), and for every compact subset K of X the set

\[
\mathrm{H} (\mathrm{K}) = \bigcup_ {x \in \mathrm{K}} \mathrm{H} (x)
\]

is compact, since it is the image of \(\mathbf{H} \times \mathbf{K}\) under the continuous mapping \((u, x) \to u(x)\) .

COROLLARY 2. Let X, Y, Z be three topological spaces such that X is Hausdorff and Y is locally compact. Then the restriction to \(\mathcal{C}(\mathbf{X} \times \mathbf{Y}; \mathbf{Z})\) of the canonical bijection \(\mathcal{F}(\mathbf{X} \times \mathbf{Y}; \mathbf{Z}) \to \mathcal{F}(\mathbf{X}; \mathcal{F}(\mathbf{Y}; \mathbf{Z}))\) (Set Theory, R, §4, no. 14) is a homeomorphism of \(\mathcal{C}_c(\mathbf{X} \times \mathbf{Y}; \mathbf{Z})\) onto \(\mathcal{C}_c(\mathbf{X}; \mathcal{C}_c(\mathbf{Y}; \mathbf{Z}))\) .

This restriction is certainly a bijection

\[
\rho : \mathcal {C} (\mathrm{X} \times \mathrm{Y}; \mathrm{Z}) \rightarrow \mathcal {C} (\mathrm{X}; \mathcal {C} _ {c} (\mathrm{Y}; \mathrm{Z}))
\]

by Theorem 3; it remains therefore to be shown that the compact-open topology on \(\mathcal{C}(\mathbf{X}\times\mathbf{Y};\mathbf{Z})\) is the inverse image under \(\rho\) of the compact-open topology on \(\mathcal{C}(\mathbf{X};\mathcal{C}_{c}(\mathbf{Y};\mathbf{Z}))\) . Since the sets \(T(\mathbf{K},\mathbf{U})\) , where \(\mathbf{K}\) is a compact subset of \(\mathbf{Y}\) and \(\mathbf{U}\) is an open subset of \(\mathbf{Z}\) , form a subbase of the topology of \(\mathcal{C}_{c}(\mathbf{Y};\mathbf{Z})\) , it follows from Remark 2 that the topology of \(\mathcal{C}_{c}(\mathbf{X};\mathcal{C}_{c}(\mathbf{Y};\mathbf{Z}))\) is generated by the sets of the form \(T(J,T(\mathbf{K},\mathbf{U}))\) , where \(\mathbf{K}\) and \(\mathbf{U}\) are as above and \(J\) is a compact subset of \(\mathbf{X}\) . Now the image of \(T(J,T(\mathbf{K},\mathbf{U}))\) under \(\rho\) is precisely \(T(J\times K,\mathbf{U})\) , and is therefore an open set; so we have shown that \(\rho\) is continuous. To prove that \(\rho\) is a homeomorphism, we note first that the sets of the form \(J\times K\) in \(\mathbf{X}\times\mathbf{Y}\) (where \(J\) is a compact subset of \(\mathbf{X}\) , and \(\mathbf{K}\) is a compact subset of \(\mathbf{Y}\) ) satisfy the condition (R) of Remark 2: for if \(\mathbf{L}\) is a compact subset of \(\mathbf{X}\times\mathbf{Y}\) and \(\mathbf{V}\) is a neighbourhood of \(\mathbf{L}\) in \(\mathbf{X}\times\mathbf{Y}\) , the projections \(M=pr_{1}(L)\) , \(N=pr_{2}(L)\) are compact, since \(\mathbf{X}\) and \(\mathbf{Y}\) are Hausdorff and \(V\cap(M\times N)\) is a neighbourhood of \(\mathbf{L}\) in the compact space \(M\times N\) , so that every point of \(\mathbf{L}\) has a neighbourhood in \(M\times N\) of the form \(J\times K\subset V\) , where \(J\subset M\) and \(K\subset N\) are compact; since \(\mathbf{L}\) can be covered by a finite number of these neighbourhoods, the assertion is proved. Therefore sets of the form \(T(J\times K;\mathbf{U})\) , where \(J\) is a compact subset of \(\mathbf{X},\mathbf{K}\) a compact subset of \(\mathbf{Y}\) and \(\mathbf{U}\) an open subset of \(\mathbf{Z}\) , generate the topology of \(\mathcal{C}_{c}(\mathbf{X}\times\mathbf{Y};\mathbf{Z})\) . But we have already seen that the image of \(T(J\times K,\mathbf{U})\) under \(\rho\) is the open set \(T(J,T(K,\mathbf{U}))\) in \(\mathcal{C}_{c}(\mathbf{X};\mathcal{C}_{c}(\mathbf{Y};\mathbf{Z}))\) ; hence \(\rho\) is a homeomorphism.

Note that if in addition Z is assumed to be uniformizable, Corollary 2 is a trivial consequence of § 1, no. 4, Proposition 2.

PROPOSITION 9. Let X, Y, Z be three topological spaces, Y being locally compact. Then the mapping \((u, v) \to v \circ u\) of \(\mathcal{C}_c(X; Y) \times \mathcal{C}_c(Y; Z)\) into \(\mathcal{C}_c(X; Z)\) is continuous.

We have to show that, for every compact subset K of X and every open subset U of Z, the set R of pairs \((u, v)\) such that \(v(u(\mathbf{K})) \subset \mathbf{U}\) is open in \(\mathcal{C}_{c}(\mathbf{X}; \mathbf{Y}) \times \mathcal{C}_{c}(\mathbf{Y}; \mathbf{Z})\) . Let \((u_{0}, v_{0}) \in R\) ; then \(u_{0}(\mathbf{K})\) is a compact subset of the locally compact space Y, contained in the open set \(\overline{v}_{0}^{1}(\mathbf{U})\) , and hence there is a compact neighbourhood L of \(u_{0}(\mathbf{K})\) contained in \(\overline{v}_{0}^{1}(\mathbf{U})\) (Chapter I, § 9, no. 7, Proposition 10). The set V of all \(u \in C_{c}(\mathbf{X}; \mathbf{Y})\) such that \(u(\mathbf{K}) \subset \mathring{\mathbf{L}}\) is a neighbourhood of \(u_{0}\) , and the set W of all \(v \in C_{c}(\mathbf{Y}; \mathbf{Z})\) such that \(v(\mathbf{L}) \subset \mathbf{U}\) is a neighbourhood of \(v_{0}\) ; furthermore, the relation \((u, v) \in \mathbf{V} \times \mathbf{W}\) implies \(v(u(\mathbf{K})) \subset \mathbf{U}\) . Hence the result.

